% Looking Away. Building 1: Moves. Parts are input below.
\documentclass[11pt,a4paper]{article}
\usepackage{amsmath,amssymb}
\usepackage[margin=2.8cm]{geometry}
\setlength{\parindent}{0pt}
\setlength{\parskip}{6pt}
% tables that wrap (array), tables that may run onto a second page (longtable),
% and the two map diagrams (tikz)
\usepackage{array,longtable}
\usepackage{tikz}
\usetikzlibrary{arrows.meta}
% number systems
\newcommand{\N}{\mathbb{N}}
\newcommand{\Z}{\mathbb{Z}}
\newcommand{\Q}{\mathbb{Q}}
\newcommand{\Zn}{\Z/n}
\newcommand{\R}{\mathbb{R}}
% build floors
\newcommand{\PP}{\mathbb{P}}   % pairs of naturals (Chapter 1)
\newcommand{\FF}{\mathbb{F}}   % selected pairs of integers (Chapter 2)
\newcommand{\Seq}{\mathcal{S}}  % sequences of rationals (Chapter 4)
\newcommand{\CC}{\mathcal{C}}    % sequences that settle (Chapter 4)
% questions
\newcommand{\sett}{\operatorname{sett}}
\newcommand{\whole}{\operatorname{whole}}
\newcommand{\mult}{\operatorname{mult}}
\newcommand{\nul}{\operatorname{null}}
\newcommand{\pos}{\operatorname{pos}}
% bags and labels
\newcommand{\bag}[1]{[#1]}
\newcommand{\lab}[1]{\textbf{#1}}
\newenvironment{lines}{\begin{list}{}{\setlength{\leftmargin}{4.2em}\setlength{\labelwidth}{3.4em}\setlength{\labelsep}{0.8em}\setlength{\itemsep}{2pt}\setlength{\topsep}{2pt}}}{\end{list}}
% links: contents entries and PDF bookmarks become clickable (keep this last)
\usepackage[colorlinks=true,linkcolor=blue,bookmarksopen=true]{hyperref}

\begin{document}
\begin{center}
{\Large\bfseries Looking Away}\\[4pt]
{\large Building 1: Moves}\\[4pt]
{\small Arrows from a floor to itself, and what counting their ``hows'' gives.}
\end{center}
\renewcommand{\contentsname}{Contents}
\tableofcontents
\bigskip

A \emph{move} is an arrow from a floor to itself that is an inside and a picture (Chapter 0, W2). Forgetting \emph{how} one thing moves to another gives a blur: Building 0. This book keeps the how. It reads Building 0 and is never read by it (the rule F7, between buildings).

Chapter numbers restart in this book; Building 0 is cited as ``Building 0, Chapter $n$''. The rules are Chapter 0 of Building 0.

\textbf{Import ledger} (Chapter 0, C8). Every line below is a Given of the form ``(Imported: source.) statement. Used by: lines.'' Nothing else in this book rests on an unearned step; classical facts no line uses sit in Remarks, marked ``reported''.
\begin{center}
\renewcommand{\arraystretch}{1.25}
\begin{tabular}{@{}l p{4.6cm} l p{4.1cm}@{}}
\textbf{Line} & \textbf{Imported} & \textbf{Used by} & \textbf{Plan to earn} \\ \hline
III.3 G5 & at most $\ell^2$ points with $\ell P = O$ & III.3 T2 & division by lines, as in III.1 P4; $\ell = 2$ is earned \\
III.3 G6 & $|N_p - p - 1| \le 2\sqrt p$ (Hasse) & III.3 T3 & the ladder $N_{p^m}$ and Amplify (Field notes: the edge, Note 2) \\
\end{tabular}
\end{center}
Earned in this pass: the ``at most one point per bag'' of the cubic (III.1 P9), by orders at a point (III.1 P1--P4); before, it rested on a deferred classical step.

Chapter kinds (Chapter 0, K1--K7) are named in each verbs line: blur, build, select, arrow, count, table, rerun. Part I and most of Part II are blur chapters; II.1 is a table, II.6, II.14, II.15 and III.3 are counts, III.2 is an arrow.

\textbf{Planned, not yet written} (\texttt{plans/ALGEBRA-PATH.md}, optional chapters): representations of finite groups; rings in several letters. The plane, the cubic and the torus are in Building 0 (Chapters 35--37); Part III here uses them.
\newpage
% Building 1, Part I. Moves on things.
% One file for the whole part. Each chapter begins with a banner and \section*.
% Building 1 reads Building 0 (cited "Building 0, Chapter n"); Building 0 never reads Building 1.

\part*{Part I. Moves on things}
\phantomsection\addcontentsline{toc}{part}{Part I. Moves on things}

%% ==================================================================
%%  Chapter 1  (rearrangements)
%% ==================================================================
\section*{Chapter 1. Rearrangements}
\phantomsection\addcontentsline{toc}{section}{\texorpdfstring{Chapter 1. Rearrangements \quad a bag of $n$ things $\to$ its moves}{Chapter 1. Rearrangements}}
\begin{center}
{\normalsize a bag of $n$ things $\to$ its moves \quad$\cdot$\quad from: Building 0, Chapters 0, 1, 3 \quad$\cdot$\quad build, select, blur}\\[4pt]
{\small Select the moves of $n$ things. Glue all swaps together and exactly two bags are left: the sign appears. Glue a move to its relabellings and the bags are the shapes of its cycles.}
\end{center}

The first chapter that keeps the \emph{how}. A move is an arrow from a bag to itself that is an inside and a picture (Chapter 0, W2). Building 0 forgot how one thing was carried to another; here the carrying is the thing.

\subsection*{Given}

\begin{lines}
\item[\lab{G1}] (Building 0, Chapter 0.) Build M1--M7; blur L1--L9; arrows W1--W4; W2: inside and picture.
\item[\lab{G2}] (Building 0, Chapter 3.) $\Z/2$ with $+$.
\item[\lab{G3}] (Building 0, Chapter 1.) Whole numbers; counting a finite bag.
\end{lines}

\subsection*{Definitions}

\begin{lines}
\item[\lab{D1}] (Build, M7; Select, M4.) Fix $n \ge 1$ and $[n] := \{1, \dots, n\}$. $S_n$: the maps $\sigma : [n] \to [n]$ that are one-to-one (then also onto). Each is a \emph{move} of $[n]$. \ $\sigma \tau := \sigma \circ \tau$ (do $\tau$ first). \ $e$: the move that does nothing.
\item[\lab{D2}] A \emph{swap} $(i\,j)$ exchanges $i$ and $j$ and keeps the rest. A \emph{cycle} $(i_1\, i_2 \dots i_k)$ sends each $i_r$ to $i_{r+1}$ and $i_k$ to $i_1$.
\item[\lab{D3}] (Inversions.) $\operatorname{inv}(\sigma)$: the number of pairs $i < j$ with $\sigma(i) > \sigma(j)$.
\item[\lab{D4}] (Relabelling.) For $\pi$ in $S_n$: $\sigma^\pi := \pi \sigma \pi^{-1}$, the move $\sigma$ with the things renamed by $\pi$.
\end{lines}

\subsection*{Propositions: the floor}

\begin{lines}
\item[\lab{P1}] (Laws, carried by W2.) $\circ$ is associative; $e$ does nothing on either side; every move has an undo $\sigma^{-1}$. Not commutative for $n \ge 3$: $(1\,2)(2\,3) = (1\,2\,3)$ while $(2\,3)(1\,2) = (1\,3\,2)$.
\item[\lab{P2}] (Count.) $S_n$ has $n!$ members.
\item[\lab{P3}] (Normal form: cycles.) Every move is a product of cycles on disjoint sets of things, in exactly one way up to the order of the cycles (follow each thing until it comes back).
\item[\lab{P4}] (Swaps are enough.) Every move is a product of swaps, by P3 and \[ (i_1 \dots i_k) = (i_1\, i_2)(i_2\, i_3) \cdots (i_{k-1}\, i_k) \quad \text{(do the right-hand swap first).} \]
\end{lines}

\subsection*{Blur 1: the sign}

On $S_n$, $n \ge 2$. \ $O$: $\circ$.

\begin{lines}
\item[\lab{B1}] (From below.) Gluings: $s \sim t$ for any two swaps $s, t$.
\item[\lab{B2}] (From above.) Question: the parity of $\operatorname{inv}(\sigma)$, in $\Z/2$.
\item[\lab{P5}] (Survives: parity respects $\circ$.) $\operatorname{inv}(\sigma\tau) \equiv \operatorname{inv}(\sigma) + \operatorname{inv}(\tau) \pmod 2$. \ (For each pair $\{i, j\}$, $\sigma\tau$ turns it over exactly when one of $\tau$ on $\{i, j\}$ and $\sigma$ on $\{\tau i, \tau j\}$ does; count the pairs mod 2.) A swap has an odd number of inversions: $(i\,j)$ with $i < j$ has $2(j - i) - 1$.
\item[\lab{P6}] (Appears.) The two blurs agree. There are exactly two bags, $A_n$ (even) and its other bag (odd). \ (B1 puts every product of an even number of swaps with $e$ and every odd product with $(1\,2)$, by P4; B2 tells them apart, by P5.) Name: $\operatorname{sgn}(\sigma) := (-1)^{\operatorname{inv}(\sigma)}$.
\item[\lab{P7}] (Dies: how many swaps.) $(1\,3) = (1\,2)(2\,3)(1\,2)$: one swap, or three. The number of swaps is not a question on moves; only its parity is.
\item[\lab{P8}] (Survives into Building 0.) $\operatorname{sgn} : S_n \to \{1, -1\}$ respects $\circ$ and $\times$: an arrow onto Building 0's $\Z/2$ (G2), written multiplicatively. By W4, $S_n \to \Z/2$ is the blur B1 followed by a matching.
\end{lines}

\subsection*{Blur 2: relabelling}

On $S_n$. \ $O$: nothing (L7: this blur respects neither $\circ$ nor its undo in general; say so).

\begin{lines}
\item[\lab{B3}] (From below.) Gluings: $\sigma \sim \sigma^\pi$ for every $\pi$.
\item[\lab{B4}] (From above.) Question: the \emph{cycle type}, the list of cycle lengths of P3 in falling order (a partition of $n$).
\item[\lab{P9}] (Appears.) The two blurs agree. \ ($\pi (i_1 \dots i_k) \pi^{-1} = (\pi i_1 \dots \pi i_k)$: relabelling renames the cycles and keeps their lengths; two moves of the same type are relabellings of each other.) Normal form: one move per partition of $n$; the bags number $p(n)$: $1, 2, 3, 5, 7, 11$ for $n = 1, \dots, 6$.
\item[\lab{P10}] (Dies: $\circ$.) $(1\,2) \sim (2\,3)$, but $(1\,2)(1\,2) = e$ while $(2\,3)(1\,2) = (1\,3\,2)$.
\item[\lab{P11}] (Survives: sign and order.) $\operatorname{sgn}(\sigma) = (-1)^{n - c}$ with $c$ the number of cycles (counting fixed things); the order of $\sigma$ (least $m \ge 1$ with $\sigma^m = e$) is the lcm of the cycle lengths. Both are read off the type.
\item[\lab{P12}] (Bag sizes.) The bag of type $1^{m_1} 2^{m_2} \cdots$ has $n! \,/\, \prod_k k^{m_k}\, m_k!$ members. \ (Count the relabellings that give the same move.)
\end{lines}

\subsection*{Theorems}

\begin{lines}
\item[\lab{T1}] (Survey: every blur of $S_n$ that respects $\circ$.) A blur of $S_n$ respecting $\circ$ is fixed by its bag of $e$, which is a union of bags of Blur 2 closed under $\circ$. For $n \le 6$ the survey is: \\
$n = 2$: two blurs. \ $n = 3$: three ($e$ only; $A_3$; everything). \ $n = 4$: four ($e$; the three double swaps with $e$; $A_4$; everything). \ $n = 5$ and $n = 6$: three ($e$; $A_n$; everything). \\
Machine-checked (all unions of types closed under $\circ$). For $n \ge 5$ the classical theorem is that $A_n$ glues to nothing smaller: $A_n$ is \emph{simple}. That is a count, proved in Part II.
\item[\lab{T2}] (The sign is the only arrow to $\Z/2$.) For $n \ge 2$, the arrows $S_n \to \{1, -1\}$ respecting $\circ$ are $\sigma \mapsto 1$ and $\operatorname{sgn}$. \ (Relabelling glues all swaps (P9), so an arrow takes one value on them; P4.)
\item[\lab{T3}] (Every move system lives here. Cayley.) For any finite bag $G$ of moves of something, closed under $\circ$ and undo, $g \mapsto (h \mapsto g h)$ is an inside of $G$ into the moves of the bag $G$, respecting $\circ$. So every finite move system is found inside some $S_n$, $n = $ its size.
\end{lines}

\subsection*{Remarks}

\begin{lines}
\item[\lab{R1}] (What the blur did.) The sign is an Appears line: ``glue all swaps'' from below meets ``parity of inversions'' from above. That the bags are two and not one is the kept question P5, and it is the whole content of the classical theorem.
\item[\lab{R2}] (Building 0 seen from here.) The blur of Building 0, Chapter 3 ($n \sim 0$ on $\Z$) is ``forget how far you moved'' for the moves $x \mapsto x + k$. Every blur in Building 0 is such a forgetting; this book keeps the moves.
\item[\lab{R3}] (Old names.) $S_n$: the symmetric group. $A_n$: the alternating group. P3: disjoint cycle decomposition. P6: the sign is well defined. B3: conjugacy; its bags are conjugacy classes. P12: class sizes. T1: the normal subgroups of $S_n$. T3: Cayley's theorem.
\item[\lab{R4}] (Machine check, \texttt{checks/b1-ch1-check.py}.) P1, P2, P4 on all of $S_n$, $n \le 6$; P5 for $n \le 5$; P6 (closure of the swap gluing); P7; P9 and P12 (types and bag sizes); P11; T1 (survey); T2.
\end{lines}

\subsection*{Exercises}

\begin{lines}
\item[\lab{E1}] Write $\sigma = (1\,4\,2)(3\,5)$ in $S_5$ as swaps, and find $\operatorname{sgn}\sigma$ and the order of $\sigma$ in two ways.
\item[\lab{E2}] Count the moves of type $2\,2\,1$ in $S_5$ by P12, and by listing.
\item[\lab{E3}] Show that in $S_n$ the even moves are exactly the products of 3-cycles ($n \ge 3$).
\item[\lab{E4}] Which move of $S_4$ has the most inversions? How many does it have?
\item[\lab{E5}] Why is ``$\sigma$ and $\tau$ commute'' not kept by Blur 2? Give a witness in $S_3$.
\end{lines}

\newpage
% Building 1, Part II. Move systems as rows.
% One file for the whole part. Each chapter begins with a banner and \section*.

\part*{Part II. Move systems as rows}
\phantomsection\addcontentsline{toc}{part}{Part II. Move systems as rows}

%% ==================================================================
%%  Chapter II.1  (the table)
%% ==================================================================
\section*{Chapter II.1. The table of move systems}
\phantomsection\addcontentsline{toc}{section}{\texorpdfstring{Chapter II.1. The table of move systems \quad move systems $\to$ classes}{Chapter II.1. The table of move systems}}
\begin{center}
{\normalsize move systems $\to$ classes \quad$\cdot$\quad from: Part I; Building 0, Chapters 1, 3, 8, 9, 14, II.1 \quad$\cdot$\quad table, blur, class}\\[4pt]
{\small The rows are now move systems. Every row passes the three laws of moves; the table is about what else they pass.}
\end{center}

Building 0, Chapter II.1 again, with move systems for rows. A \emph{move system} is a bag $G$ with $\circ$, a do-nothing $e$ and undos, where $\circ$ is associative: exactly the laws Chapter 1, P1 found for the moves of a bag. By Chapter 1, T3 every finite one is a bag of moves.

\subsection*{Given}

\begin{lines}
\item[\lab{G1}] (Rows.) $\Z$ with $+$ (Building 0, Ch.\ 1); $\Z/6$ (Ch.\ 3); the circle $\mathbb{T}$ (Ch.\ 8); $\Q^\times$, the non-zero rationals with $\times$ (Ch.\ 2); $(\Z/8)^\times = \{1, 3, 5, 7\}$ with $\times$ (Ch.\ 3); $S_3$ (Chapter 1); $D_4$, the moves of $S_4$ that carry a square's corners $1, 2, 3, 4$ (in order round the square) to corners of the square; $GL_2(\R)$, the straight maps of $\R^2$ with $\det \ne 0$ (Building 0, Ch.\ 14); $F_2$ (Chapter 2); $\Z^2$ (Chapter 2, T3).
\item[\lab{G2}] (Facts.) The lines of those chapters, read as yes/no facts.
\end{lines}

\subsection*{Tests}

\begin{lines}
\item[\lab{D1}] $\mathsf{comm}$: $gh = hg$ for all $g, h$.
\item[\lab{D2}] $\mathsf{fin}$: finitely many members.
\item[\lab{D3}] $\mathsf{cyc}$: some $g$ has every member as a power $g^k$, $k$ in $\Z$.
\item[\lab{D4}] $\mathsf{tf}$: no $g \ne e$ has $g^k = e$ for any $k \ge 1$.
\item[\lab{D5}] $\mathsf{div}$: for every $g$ and every $n \ge 1$ there is $h$ with $h^n = g$.
\end{lines}

\subsection*{The table}

\begin{center}
\renewcommand{\arraystretch}{1.25}
\begin{tabular}{l|ccccc}
 & $\mathsf{comm}$ & $\mathsf{fin}$ & $\mathsf{cyc}$ & $\mathsf{tf}$ & $\mathsf{div}$ \\ \hline
$\Z$ & $\checkmark$ & $\cdot$ & $\checkmark$ & $\checkmark$ & $\cdot$ \\
$\Z/6$ & $\checkmark$ & $\checkmark$ & $\checkmark$ & $\cdot$ & $\cdot$ \\
$\mathbb{T}$ & $\checkmark$ & $\cdot$ & $\cdot$ & $\cdot$ & $\checkmark$ \\
$\Q^\times$ & $\checkmark$ & $\cdot$ & $\cdot$ & $\cdot$ & $\cdot$ \\
$(\Z/8)^\times$ & $\checkmark$ & $\checkmark$ & $\cdot$ & $\cdot$ & $\cdot$ \\
$S_3$ & $\cdot$ & $\checkmark$ & $\cdot$ & $\cdot$ & $\cdot$ \\
$D_4$ & $\cdot$ & $\checkmark$ & $\cdot$ & $\cdot$ & $\cdot$ \\
$GL_2(\R)$ & $\cdot$ & $\cdot$ & $\cdot$ & $\cdot$ & $\cdot$ \\
$F_2$ & $\cdot$ & $\cdot$ & $\cdot$ & $\checkmark$ & $\cdot$ \\
$\Z^2$ & $\checkmark$ & $\cdot$ & $\cdot$ & $\checkmark$ & $\cdot$ \\
\end{tabular}
\end{center}
Reasons for the less obvious entries: $\mathbb{T}$ is not $\mathsf{cyc}$ (if $a$ turns a rational amount its powers are finitely many; if not, $a/2$ is no power of $a$) and is $\mathsf{div}$ (Building 0, Chapter 8, T2). $\Q^\times$: $-1$ has $(-1)^2 = 1$; $2$ has no square root. $(\Z/8)^\times$: every member squares to $1$. $GL_2(\R)$: $-I$ squares to $I$; $\operatorname{diag}(-1, 1)$ has no square root ($\det$ of a square is $\ge 0$). $F_2$: $s$ has no square root (exponent sums, Chapter 2, P5); no power of a reduced word $\ne 1$ is $1$.

\subsection*{Blur, and what appears}

\begin{lines}
\item[\lab{B1}] (Building 0, II.1, B1.) For a set $U$ of tests, glue two rows with the same answers on $U$.
\item[\lab{P1}] (Not sharp.) With all five tests, $S_3$ and $D_4$ share a bag. A separating test is the size: $6$ against $8$. Every other pair of rows is told apart.
\item[\lab{P2}] (Classes.) As $T$ runs over the sets of tests, $\operatorname{class}(T)$ takes $11$ values. Old names: $\operatorname{class}(\mathsf{comm})$, \emph{abelian}; $\operatorname{class}(\mathsf{cyc})$, \emph{cyclic} ($\Z$, $\Z/6$); $\operatorname{class}(\mathsf{tf})$, \emph{torsion-free} ($\Z$, $F_2$, $\Z^2$).
\item[\lab{P3}] (Read off the table.) $\mathsf{cyc} \Rightarrow \mathsf{comm}$; \ $\mathsf{div} \Rightarrow \mathsf{comm}$. Never together: $\mathsf{fin}$ and $\mathsf{tf}$; $\mathsf{fin}$ and $\mathsf{div}$; $\mathsf{cyc}$ and $\mathsf{div}$; $\mathsf{tf}$ and $\mathsf{div}$.
\end{lines}

\subsection*{Theorems}

For every move system with more than one member, not only the ten rows.
\begin{lines}
\item[\lab{T1}] $\mathsf{cyc} \Rightarrow \mathsf{comm}$. \ ($g^j g^k = g^{j+k} = g^k g^j$.)
\item[\lab{T2}] $\mathsf{fin}$ and $\mathsf{tf}$ never together. \ (The powers of $g$ repeat, $g^i = g^j$ with $i < j$, so $g^{j-i} = e$.)
\item[\lab{T3}] $\mathsf{fin}$ and $\mathsf{div}$ never together. \ (By T2's reason each $g$ has $g^m = e$ for some $1 \le m \le |G|$; so with $n = |G|!$ every $n$-th power is $e$, and $g \ne e$ has no $n$-th root.)
\item[\lab{T4}] $\mathsf{cyc}$ and $\mathsf{div}$ never together. \ (Finite: T3. Infinite: it matches $\Z$, and $2h = 1$ has no answer there.)
\item[\lab{T5}] (Matching rows.) $\Z/6$ matches $\Z/2 \times \Z/3$, and $(\Z/8)^\times$ matches $\Z/2 \times \Z/2$; $\Z/6$ and $(\Z/8)^\times$ differ already by size, and $\Z/4$ and $(\Z/8)^\times$ by $\mathsf{cyc}$.
\end{lines}

\subsection*{Remarks}

\begin{lines}
\item[\lab{R1}] (Accidents of a small table; Building 0, II.1, R1.) Two lines of P3 die when the right row is added. $\mathsf{div} \Rightarrow \mathsf{comm}$ dies on the rotations of space: every rotation has $n$-th roots (turn by a $1/n$ of the angle about the same axis), and two rotations about different axes need not commute. ``$\mathsf{tf}$ and $\mathsf{div}$ never together'' dies on $\Q$ with $+$. The other four survive every row: T1--T4.
\item[\lab{R2}] (A conjecture from the table, F6.) Every abelian row with $\mathsf{fin}$ in the table matches a product of clocks $\Z/n_1 \times \dots \times \Z/n_k$. This is true for every finite abelian move system; it is proved in the chapter on modules, from a normal form of matrices.
\item[\lab{R3}] (Old names.) Move system: group. $\mathsf{comm}$: abelian. $\mathsf{tf}$: torsion-free. $\mathsf{div}$: divisible. $D_4$: the dihedral group of order $8$. T2: finite groups are torsion groups.
\item[\lab{R4}] (Machine check, \texttt{checks/b1-p2-ch1-check.py}.) Every finite entry by brute force; the $11$ classes (\texttt{b1-p2-ch1-classes.py}); P1; T5 (the two matchings checked).
\end{lines}

\subsection*{Exercises}

\begin{lines}
\item[\lab{E1}] List the members of $D_4$ as moves of the four corners, and show $D_4$ is not $\mathsf{comm}$.
\item[\lab{E2}] Add the row $\Q$ with $+$. How many classes now?
\item[\lab{E3}] Find a test that separates $S_3$ from $D_4$ without counting members.
\item[\lab{E4}] Show $\Z^2$ is not $\mathsf{cyc}$.
\end{lines}

\newpage

%% ==================================================================
%%  Chapter II.2  (cosets)
%% ==================================================================
\section*{Chapter II.2. Blurs that respect one side}
\phantomsection\addcontentsline{toc}{section}{\texorpdfstring{Chapter II.2. Blurs that respect one side \quad a move system $\to$ its cosets}{Chapter II.2. Blurs that respect one side}}
\begin{center}
{\normalsize a move system $\to$ its cosets \quad$\cdot$\quad from: Chapter II.1; Building 0, Chapters 3, II.2 \quad$\cdot$\quad blur, count}\\[4pt]
{\small Ask a blur to respect moving on the left only. Its bags are copies of the bag of $e$, all of one size. Counting them gives Lagrange, and Fermat's little theorem falls out.}
\end{center}

Building 0, II.2 asked which blurs of a ring respect $+$ and $\times$. Here the question is asked of a move system with less: respect $g \circ$ on the left, and nothing on the right. Rule L7 (say what the blur respects) is the whole story of this chapter and the next.

\subsection*{Given}

\begin{lines}
\item[\lab{G1}] (Chapter II.1.) Move systems; the rows.
\item[\lab{G2}] (Building 0, Chapter 3.) $\Z/n$; $(\Z/n)^\times$, the bags with an inverse for $\times$; Chapter 3, T5 there.
\end{lines}

\subsection*{The blur}

On a move system $G$. \ $O$: for each $g$, the move $x \mapsto g x$ (left only).

\begin{lines}
\item[\lab{B1}] (From below.) For a bag $H \subseteq G$: glue $h \sim e$ for each $h$ in $H$.
\item[\lab{B2}] (From above.) For a sub move system $K$ (closed under $\circ$ and undo, holding $e$): the question $x \mapsto$ the bag $xK := \{x k : k \in K\}$.
\end{lines}

\subsection*{Propositions}

\begin{lines}
\item[\lab{P1}] (Survey: every such blur.) A blur of $G$ respecting left moves is fixed by its bag of $e$, which is a sub move system $K$; and $x \sim y \iff x^{-1} y \in K$. Conversely every sub move system gives one. So these blurs and the sub move systems match.
\item[\lab{P2}] (Appears.) B1 for $H$ equals B2 for $K = $ the smallest sub move system holding $H$. Its bags are the $x K$: the \emph{left cosets}.
\item[\lab{P3}] (Survives: moving a bag.) $x \mapsto g x$ carries the bag $yK$ onto the bag $gyK$, one-to-one. \ (Its undo is $x \mapsto g^{-1} x$.) So all bags have the same size as $K$.
\item[\lab{P4}] (Dies: the right side.) In $S_3$ with $K = \{e, (1\,2)\}$: $e \sim (1\,2)$, but $e (1\,3) = (1\,3)$ and $(1\,2)(1\,3) = (1\,3\,2)$ are in different bags, since $(1\,3)^{-1}(1\,3\,2) = (2\,3) \notin K$.
\item[\lab{P5}] (Normal form.) None in general: no rule picks one member of each $xK$. For $\Z$ and $K = n\Z$ the normal forms $0, \dots, n - 1$ of Building 0, Chapter 3 are such a rule.
\end{lines}

\subsection*{Theorems}

\begin{lines}
\item[\lab{T1}] (Counting bags: Lagrange. A count, not a Survives line.) If $G$ is finite, $|G| = |K| \cdot (\text{number of bags})$. \ (P3: the bags are disjoint and each has $|K|$ members.) So $|K|$ divides $|G|$.
\item[\lab{T2}] (The order of a move divides the size.) For $g$ in a finite $G$, the least $m \ge 1$ with $g^m = e$ divides $|G|$. \ ($K = \{e, g, \dots, g^{m-1}\}$ has $m$ members: T1.)
\item[\lab{T3}] ($g^{|G|} = e$.) From T2.
\item[\lab{T4}] (Building 0, Chapter 3, T5, recovered.) For $a$ with $\gcd(a, n) = 1$: $a^{\varphi(n)} \equiv 1 \pmod n$, where $\varphi(n) = |(\Z/n)^\times|$. For $n = p$ prime: $a^{p-1} \equiv 1$. \ (T3 in $(\Z/n)^\times$.)
\item[\lab{T5}] (Dies: ``every divisor is a size''.) $A_4$ has $12$ members and no sub move system with $6$. \ (Machine-checked; classical.) T1 is ``at most these sizes'', never ``these sizes occur''.
\item[\lab{T6}] (Prime size.) If $|G| = p$ is prime, $G$ matches $\Z/p$. \ (Any $g \ne e$ has order dividing $p$, so $p$: its powers are all of $G$.)
\end{lines}

\subsection*{Remarks}

\begin{lines}
\item[\lab{R1}] (What the blur did.) The cosets are the bags of the one blur that respects left moves; their equal size is a Survives line (P3); Lagrange is a count of bags (T1). The method's own split: blur keeps (P3), counting gives existence and sizes.
\item[\lab{R2}] (Respecting both sides.) If the blur also respects right moves, $K$ must be kept by relabelling, $gKg^{-1} = K$, and the bags themselves form a move system. That is Chapter II.3.
\item[\lab{R3}] (Old names.) Left cosets; index $= $ number of bags; T1: Lagrange's theorem; T4: Euler's theorem, Fermat's little theorem; T5: the converse of Lagrange fails ($A_4$).
\item[\lab{R4}] (Machine check, \texttt{checks/b1-p2-ch2-check.py}.) P1 (for every bag $K \ni e$ of $S_3$ and $D_4$: ``$x^{-1}y \in K$'' is a blur exactly when $K$ is a sub move system); P3; P4; T1--T3 on all sub move systems of $S_4$; T4 for $n \le 200$; T5.
\end{lines}

\subsection*{Exercises}

\begin{lines}
\item[\lab{E1}] List the left cosets of $K = \{e, (1\,2)\}$ in $S_3$, and then the right cosets $Kx$. Are they the same bags?
\item[\lab{E2}] Find $3^{100} \bmod 7$ with T4.
\item[\lab{E3}] Which sizes of sub move systems occur in $S_4$? Compare with the divisors of $24$.
\item[\lab{E4}] Show that a move system of size $4$ is $\mathsf{comm}$. (Each $g \ne e$ has order $2$ or $4$.)
\end{lines}

\newpage

%% ==================================================================
%%  Chapter II.3  (normal bags)
%% ==================================================================
\section*{Chapter II.3. Blurs that respect both sides}
\phantomsection\addcontentsline{toc}{section}{\texorpdfstring{Chapter II.3. Blurs that respect both sides \quad a move system $\to$ its quotients}{Chapter II.3. Blurs that respect both sides}}
\begin{center}
{\normalsize a move system $\to$ its quotients \quad$\cdot$\quad from: Chapters 1, 2, II.2 \quad$\cdot$\quad blur}\\[4pt]
{\small Ask the blur to respect moving on both sides. The bag of $e$ must now be kept by relabelling, and the bags form a move system of their own. Every arrow is such a blur followed by an inside.}
\end{center}

The same gluing as Chapter II.2, with one more operation in $O$. Building 0, Chapter 3 is the case $G = \Z$, where left and right are the same and every sub move system will do.

\subsection*{Given}

\begin{lines}
\item[\lab{G1}] (Chapter II.2.) Left-respecting blurs and cosets; T1 (Lagrange).
\item[\lab{G2}] (Chapter 1.) Relabelling $\sigma^\pi$; the bags of relabelling (types); T1 (survey for $S_n$, $n \le 6$).
\item[\lab{G3}] (Chapter 2.) Presentations; P5 (exponent sums).
\end{lines}

\subsection*{The blur}

On a move system $G$. \ $O$: $x \mapsto gx$ and $x \mapsto xg$ for each $g$.

\begin{lines}
\item[\lab{B1}] (From below.) For a bag $H$: glue $h \sim e$ for each $h$ in $H$.
\item[\lab{B2}] (From above.) For a sub move system $N$ with $gNg^{-1} = N$ for every $g$ (\emph{kept by relabelling}): the question $x \mapsto xN$.
\end{lines}

\subsection*{Propositions}

\begin{lines}
\item[\lab{P1}] (Survey.) A blur respecting both sides is a left-respecting blur (Chapter II.2, P1) whose bag of $e$, $N$, is kept by relabelling: $n \sim e$ gives $g n g^{-1} \sim g g^{-1} = e$. Conversely each such $N$ gives one. Then $xN = Nx$: left and right bags agree.
\item[\lab{P2}] (Appears.) B1 for $H$ equals B2 for the smallest $N$ kept by relabelling that holds $H$ (Chapter 2, T2).
\item[\lab{P3}] (Survives: $\circ$ on bags.) $xN \cdot yN := xyN$ does not depend on the members chosen. The bags form a move system $G/N$, and $x \mapsto xN$ is a picture respecting $\circ$.
\item[\lab{P4}] (Dies: ``kept by relabelling'' does not chain.) In $D_4$ (corners $1, 2, 3, 4$ round a square), let $s$ be the flip $(2\,4)$ and $V = \{e, s, (1\,3), (1\,3)(2\,4)\}$. Then $\{e, s\}$ is kept by relabelling in $V$ ($V$ is $\mathsf{comm}$), and $V$ is kept in $D_4$ (it has half the members). But $\{e, s\}$ is not kept in $D_4$: the quarter turn $r = (1\,2\,3\,4)$ gives $r s r^{-1} = (1\,3)$.
\item[\lab{P5}] (Half is always kept.) If $N$ has half the members of $G$, it is kept by relabelling: the two left bags and the two right bags are both $N$ and the rest.
\end{lines}

\subsection*{Theorems}

\begin{lines}
\item[\lab{T1}] (Every arrow is a blur, then an inside. W4.) For an arrow $f : G \to G'$ respecting $\circ$, the blur ``$f(x) = f(y)$'' respects both sides; its bag of $e$ is $N = \{x : f(x) = e\}$; and $G/N \to G'$, $xN \mapsto f(x)$, is an inside. So $G/N$ matches the picture of $f$. \\
Example: $\operatorname{sgn} : S_n \to \{1, -1\}$ gives $S_n / A_n$ matching $\Z/2$ (Chapter 1, P8).
\item[\lab{T2}] (Blurs of a blur.) The two-sided blurs of $G/N$ match the two-sided blurs of $G$ that are blurrier than $N$'s, by $M/N \leftrightarrow M$; and $(G/N)/(M/N)$ matches $G/M$. Blurring twice is blurring once, with the coarser bag of $e$.
\item[\lab{T3}] (A sub and a blur.) For a sub move system $K$ and $N$ kept by relabelling: $KN$ is a sub move system, $K \cap N$ is kept by relabelling in $K$, and $KN/N$ matches $K/(K \cap N)$. \ (T1 for $K \to G/N$.)
\item[\lab{T4}] (Appears: the blur to $\mathsf{comm}$.) Glue $xy \sim yx$ for all $x, y$, respecting both sides. The bags form a $\mathsf{comm}$ move system $G^{\mathrm{ab}}$, and every arrow from $G$ to a $\mathsf{comm}$ move system passes through it, in exactly one way. From above: the questions are all arrows to $\mathsf{comm}$ systems. \\
$S_n^{\mathrm{ab}}$ matches $\Z/2$, by $\operatorname{sgn}$ ($n \ge 2$). \ $F_2^{\mathrm{ab}}$ matches $\Z^2$, by the exponent sums (Chapter 2, P5). \ $D_4^{\mathrm{ab}}$ matches $\Z/2 \times \Z/2$.
\item[\lab{T5}] (Simple: only two such blurs. A count.) $G$ is \emph{simple} if its only two-sided blurs are ``glue nothing'' and ``glue everything''. $\Z/p$ is simple (Chapter II.2, T1). $A_5$ is simple: its relabelling bags (inside $A_5$) have sizes $1, 15, 20, 12, 12$; a bag of $e$ kept by relabelling is a union of these holding the $1$, and its size divides $60$ (Lagrange); the only such unions have $1$ or $60$ members.
\end{lines}

\subsection*{Remarks}

\begin{lines}
\item[\lab{R1}] (L7 is the whole story.) The same gluing $h \sim e$, with left moves in $O$, gives cosets; with both sides, quotients. Nothing else changed.
\item[\lab{R2}] (Simple pieces.) A ladder of blurs $G \to G/N_1 \to \dots \to \{e\}$ in which each rung's blur is simple lists the \emph{atoms} of $G$ (the atoms ladder of the project, rung 2: the list survives, the gluing does not). $A_5$ is the first atom that is not a clock.
\item[\lab{R3}] (Old names.) Kept by relabelling: normal subgroup. $G/N$: quotient group. T1: the first isomorphism theorem. T2: the correspondence and third isomorphism theorems. T3: the second isomorphism theorem. T4: abelianisation, $G/[G, G]$. T5: simple groups; $A_5$ simple. R2: Jordan--H\"older.
\item[\lab{R4}] (Machine check, \texttt{checks/b1-p2-ch3-check.py}.) P1 on all bags of $S_3$, $D_4$ (and kept-by-relabelling sub systems of $S_4$); P3; P4; P5 on $S_4$; T1 for $\operatorname{sgn}$ and $\det$ on $GL_2(\Z/3)$; T4 sizes for $S_3$, $S_4$, $D_4$; T5 (class sizes of $A_5$, all unions).
\end{lines}

\subsection*{Exercises}

\begin{lines}
\item[\lab{E1}] List the sub move systems of $S_3$ kept by relabelling.
\item[\lab{E2}] Show that every sub move system of a $\mathsf{comm}$ system is kept by relabelling. Which line of Building 0, Chapter 3 is this?
\item[\lab{E3}] Find $S_3^{\mathrm{ab}}$ by gluing $(1\,2)(2\,3) \sim (2\,3)(1\,2)$ and closing up.
\item[\lab{E4}] Why can $A_5$ have no arrow onto $\Z/2$?
\end{lines}

\newpage

%% ==================================================================
%%  Chapter II.4  (modules and the Smith form)
%% ==================================================================
\section*{Chapter II.4. Moves on matrices: the Smith form}
\phantomsection\addcontentsline{toc}{section}{\texorpdfstring{Chapter II.4. Moves on matrices: the Smith form \quad relation matrices $\to$ their normal form}{Chapter II.4. Moves on matrices: the Smith form}}
\begin{center}
{\normalsize relation matrices $\to$ their normal form \quad$\cdot$\quad from: Chapters II.1, II.3; Building 0, Chapters 18, II.2, II.5 \quad$\cdot$\quad blur, build}\\[4pt]
{\small Glue a matrix to every matrix reached by moving its rows and columns. The kept questions are the gcds of the minors; the normal form is a diagonal in which each entry goes into the next. Every finite $\mathsf{comm}$ move system is then a product of clocks.}
\end{center}

The first of five chapters with one template: a move system acts on matrices, the blur forgets the moves, kept questions are found, and a normal form shows they are all. Here the moves are on both sides, $A \mapsto PAQ$.

\subsection*{Given}

\begin{lines}
\item[\lab{G1}] (Building 0, Chapter II.5.) $K[x]$; division with remainder by any $d \ne 0$ when $K$ is a field, and in $\Z$; $\det$ over any ring; P10 (equations are carried); P11 (inverse iff $\det$ is a unit); T2: $\Z$ and $K[x]$ pass the same lines.
\item[\lab{G2}] (Building 0, Chapter 18.) $\gcd$, and ``goes into'' as a one-way blur.
\item[\lab{G3}] (Building 0, Chapter II.2.) Blurs of a ring by ideals; quotients $R/I$.
\item[\lab{G4}] (Chapter II.3.) Quotients of move systems; T1.
\end{lines}

\subsection*{Definitions}

\begin{lines}
\item[\lab{D1}] $R$: $\Z$, or $K[x]$ with $K$ a field. Its units: $\pm 1$, or the non-zero constants.
\item[\lab{D2}] An \emph{$R$-module}: a bag $M$ with $+$ (a $\mathsf{comm}$ move system) and scaling by members of $R$, with the laws of Building 0, Chapter 6, P2. $R^n$ is one; so is every $\mathsf{comm}$ move system, with $R = \Z$ ($k \cdot m := m + \dots + m$).
\item[\lab{D3}] (Build, then blur.) For an $m \times n$ matrix $A$ over $R$: $\operatorname{coker} A := R^n$ blurred by gluing each row of $A$ to $0$, respecting $+$ and scaling. Every finite $\mathsf{comm}$ move system $M$ is a $\operatorname{coker} A$: one generator $x_g$ for each member $g$, one relation $x_g + x_h - x_{g+h}$ for each pair (Chapter 2, T1 and Chapter II.3, T1, read with $+$).
\item[\lab{D4}] $d_k(A)$: the $\gcd$ of all $k \times k$ minors of $A$ (made monic, or $\ge 0$); $d_0 := 1$.
\end{lines}

\subsection*{The blur}

On $m \times n$ matrices over $R$. \ $O$: none beyond the moves themselves.

\begin{lines}
\item[\lab{B1}] (From below.) Gluings: $A \sim PAQ$ for $P$, $Q$ square with $\det$ a unit (G1, P11). Equivalently, by row and column moves: swap two rows; add $r$ times one row to another; scale a row by a unit; and the same for columns.
\item[\lab{B2}] (From above.) Questions: $d_1(A), d_2(A), \dots$.
\end{lines}

\subsection*{Propositions}

\begin{lines}
\item[\lab{P1}] (Survives: the $d_k$.) Each $k \times k$ minor of $PA$ is an $R$-combination of $k \times k$ minors of $A$ (an equation in the entries, carried to every $R$: G1, P10). So $d_k(A)$ goes into $d_k(PA)$, and back again by $P^{-1}$: $d_k$ is kept. The same for $AQ$.
\item[\lab{P2}] (Survives: the module.) $\operatorname{coker}(PAQ)$ matches $\operatorname{coker} A$. \ ($P$ only remixes the relations; $Q$ renames the generators.)
\item[\lab{P3}] (Dies: the entries.) $\begin{pmatrix} 2 & 0 \\ 0 & 3 \end{pmatrix} \sim \begin{pmatrix} 1 & 0 \\ 0 & 6 \end{pmatrix}$ over $\Z$: $d_1 = 1$, $d_2 = 6$ for both.
\item[\lab{P4}] (Normal form; at least one, by Euclid.) Every $A$ is glued to a diagonal \[ \operatorname{diag}(e_1, e_2, \dots, e_r, 0, \dots) \] with $e_1 \mid e_2 \mid \dots \mid e_r$, all $\ne 0$. \ (Move a smallest non-zero entry to the corner; divide with remainder along its row and column; if a remainder is left, it is smaller, and repeat; sizes or degrees cannot fall forever (Building 0, Chapter 1, G8). Then clear the row and column, and make the corner go into every other entry by adding a row.)
\item[\lab{P5}] (At most one.) $e_1 e_2 \cdots e_k = d_k(A)$ up to a unit, so $e_k = d_k / d_{k-1}$: fixed by the kept questions.
\item[\lab{P6}] (Appears.) B1 and B2 agree: two matrices are glued exactly when all their $d_k$ agree. \ (P1 one way; P4 and P5 the other.)
\item[\lab{P7}] (Dies over $\Z[x]$.) The $1 \times 2$ matrix $(2 \;\; x)$ over $\Z[x]$ has $d_1 = 1$, but is glued to no $(e \;\; 0)$: the multiples of $2$ and $x$ together are not the multiples of one member (the constant terms of $2a + xb$ are even). Division with remainder was the line that P4 used, and $\Z[x]$ fails it.
\end{lines}

\subsection*{Theorems}

\begin{lines}
\item[\lab{T1}] (Every module given by finitely many generators and relations.) It matches \[ R^{s} \times R/(e_1) \times \dots \times R/(e_r), \qquad e_1 \mid e_2 \mid \dots \mid e_r, \] for exactly one $s$ and one such list of $e_i$ that are not units (up to units). \ (P2, P4, P5.)
\item[\lab{T2}] (Chapter II.1, R2, proved.) Every finite $\mathsf{comm}$ move system matches a product of clocks $\Z/e_1 \times \dots \times \Z/e_r$ with $e_1 \mid \dots \mid e_r$, in one way only. \ ($R = \Z$, $s = 0$.)
\item[\lab{T3}] (Counting them.) The $\mathsf{comm}$ move systems with $n$ members, up to matching, number $\prod_p \mathrm{p}(a_p)$, where $n = \prod p^{a_p}$ and $\mathrm{p}(a)$ is the number of partitions of $a$. For $n = 8$: $3$ ($\Z/8$, $\Z/2 \times \Z/4$, $\Z/2 \times \Z/2 \times \Z/2$). For $n = 72$: $6$. \ (Split each $e_i$ by primes: Building 0, Chapter 3, T6.)
\item[\lab{T4}] (Two clocks, again.) $\Z/m \times \Z/n$ matches $\Z/mn$ exactly when $\gcd(m, n) = 1$: the Smith form of $\operatorname{diag}(m, n)$ is $\operatorname{diag}(\gcd, \operatorname{lcm})$.
\end{lines}

\subsection*{Remarks}

\begin{lines}
\item[\lab{R1}] (The template.) A move system acting on matrices; the blur ``reachable by a move''; kept questions; a normal form proving they are all. Next in Building 1: equivalence over a field (rank), similarity (Jordan), forms (Sylvester), and the cycle type of Chapter 1 was the first case.
\item[\lab{R2}] (Where existence came from.) Once: the Euclid steps of P4, a construction. Uniqueness came from kept questions (P5).
\item[\lab{R3}] (Old names.) $R$-module; $\operatorname{coker} A$: a presented module. B1: equivalence of matrices. $d_k$: determinantal divisors. P4: the Smith normal form; $e_i$: invariant factors. T1: the structure theorem for finitely generated modules over a principal ideal domain (here Euclidean). T2: the fundamental theorem of finite abelian groups. P7: $\Z[x]$ is not a principal ideal domain.
\item[\lab{R4}] (Machine check, \texttt{checks/b1-p2-ch4-check.py}.) P1 ($d_k$ kept under $200$ random moves); P4--P6 against a Smith form routine on random integer matrices; P3; T3 for $n \le 200$ by listing diagonal forms; T4.
\end{lines}

\subsection*{Exercises}

\begin{lines}
\item[\lab{E1}] Find the Smith form of $\begin{pmatrix} 2 & 4 & 4 \\ -6 & 6 & 12 \\ 10 & -4 & -16 \end{pmatrix}$ over $\Z$, and say which $\mathsf{comm}$ move system it presents.
\item[\lab{E2}] List the $\mathsf{comm}$ move systems with $36$ members.
\item[\lab{E3}] Over $\mathbb{Q}[x]$, find the Smith form of $\begin{pmatrix} x & 1 \\ 0 & x \end{pmatrix}$.
\item[\lab{E4}] Show $\Z/4 \not\cong \Z/2 \times \Z/2$ by P5, and by $\mathsf{cyc}$.
\end{lines}

\newpage
% Building 1, Part III. Curves.
% One file for the whole part. Each chapter begins with a banner and \section*.

\part*{Part III. A cubic and its moves}
\phantomsection\addcontentsline{toc}{part}{Part III. A cubic and its moves}

Building 0, Chapters 35--37 drew a cubic in the plane, found the functions on it, and met the same curve over $\mathbb{C}$ as a torus. This part glues its points by lines and finds a $\mathsf{comm}$ move system, then reads it over $\mathbb{C}$ and over $\Z/p$.

%% ==================================================================
%%  Chapter III.1  (the group law as a blur)
%% ==================================================================
\section*{Chapter III.1. Points glued by lines}
\phantomsection\addcontentsline{toc}{section}{\texorpdfstring{Chapter III.1. Points glued by lines \quad sums of points $\to$ points}{Chapter III.1. Points glued by lines}}
\begin{center}
{\normalsize sums of points $\to$ points \quad$\cdot$\quad from: Chapters 2, II.1, II.3; Building 0, Chapters 35, 36 \quad$\cdot$\quad build, blur}\\[4pt]
{\small Add points as letters, and glue the three points of each line to three copies of $O$. The bags form a move system for free, and each bag of size $0$ holds exactly one point.}
\end{center}

The usual road defines a sum of two points by a drawing and then spends pages proving it is associative. Here the sum is a blur of a free move system, so associativity is never in doubt; the work moves to ``at most one point per bag''.

\subsection*{Given}

\begin{lines}
\item[\lab{G1}] (Building 0, Chapter 36.) $k$ with $2, 3 \ne 0$; the cubic $C$: $y^2 = x^3 + ax + b$ with $\Delta \ne 0$, and $O$. P1 (the edge meets $C$ only at $O$, three times); P2 (a line meets $C$ in at most $3$, counted); P3 (the mirror); T1 (Build: two points of $C(k)$ on a line give a third, and the counts add to exactly $3$); P10 (no function has weight $1$); P11.
\item[\lab{G2}] (Chapter 2, T1; Chapter II.1.) The free $\mathsf{comm}$ move system on a bag: finite sums $\sum n_P\,P$, $n_P$ in $\Z$.
\item[\lab{G3}] (Chapter II.3.) The bags of a blur of a move system that respects both sides form a move system; for a $\mathsf{comm}$ one, respecting $+$ is enough.
\end{lines}

\subsection*{Definitions}

\begin{lines}
\item[\lab{D1}] (Build.) $\Z[C]$: the free $\mathsf{comm}$ move system on the points of $C(k)$ (G2). \\ $\deg \sum n_P P := \sum n_P$.
\item[\lab{D2}] For a line $L$, $L \cdot C := $ its meeting points with their counts, as a member of $\Z[C]$: $P + Q + R$, or $2P + R$ for a tangent, or $3O$ for the edge.
\end{lines}

\subsection*{The blur}

On $\Z[C]$. \ $O$: $+$.

\begin{lines}
\item[\lab{B1}] (From below.) Gluing: $L \cdot C \sim 3\,O$ for every line $L$ that meets $C$ at points of $C(k)$ (G1, T1: all three are then in $C(k)$). Respects $+$. \ Write $\operatorname{Pic}(C)$ for the bags.
\item[\lab{B2}] (From above.) Questions: the degree; and, for each $D$ of degree $0$, ``is $D$ glued to $0$?'', read through functions: $D \sim 0$ exactly when $D$ is the list of zeros minus poles of a product of powers of lines (P2).
\end{lines}

\subsection*{Propositions}

\begin{lines}
\item[\lab{P1}] (Free: a move system.) $\operatorname{Pic}(C)$ is a $\mathsf{comm}$ move system (G3). Its law is associative because $+$ on $\Z[C]$ is. Nothing was drawn.
\item[\lab{P2}] (What a gluing is.) The affine line $\ell = \alpha x + \beta y + \gamma$ is a function on $C$ (Building 0, Chapter 36) whose zeros, counted, are the affine points of $L \cdot C$, and whose only pole is at $O$, of order $w(\ell) = 3$ or $2$ (P11 there). So $L \cdot C - 3O$ is ``zeros minus poles'' of $\ell$, and $D \sim 0$ exactly when $D$ is zeros minus poles of $\prod \ell_i^{n_i}$.
\item[\lab{P3}] (Survives: the degree.) Each gluing has degree $3 - 3 = 0$. So $\deg$ is kept, and $\operatorname{Pic}^0(C)$, the bags of degree $0$, is a sub move system.
\item[\lab{P4}] (Two moves on points.) For $P$, $Q$ in $C(k)$ let $P * Q$ be the third point on their line (G1, T1). Then \[ P + Q + (P * Q) \sim 3O, \qquad R + \bar R \sim 2O, \] the second by the upright line through $R$ and its mirror $\bar R$ (G1, P3). So $P + Q \sim O + \overline{P * Q}$: two points are glued to one point plus $O$.
\end{lines}

\subsection*{Theorems}

\begin{lines}
\item[\lab{T1}] (At least one: normal form.) Every $D$ of degree $0$ is glued to $P - O$ for some $P$ in $C(k)$. \ (Write $D = \sum (P_i - O) - \sum (Q_j - O)$. $-(Q - O) \sim \bar Q - O$ by P4. Then $(P - O) + (Q - O) \sim \overline{P * Q} - O$ by P4, and repeat: each step is a Build of Building 0, Chapter 36, T1.)
\item[\lab{T2}] (At most one.) $P - O \sim Q - O$ only when $P = Q$. \ (Put $R := \overline{P * \bar Q}$; then $P - Q \sim R - O$, so $R - O$ is zeros minus poles of some $f = \prod \ell_i^{n_i}$ (P2). If $R \ne O$, $f$ has no pole on the affine part, so $f$ lies in $k[C]$ (classical, proof deferred: a function on $C$ with no affine pole is a polynomial function); its only pole is a simple one at $O$, so $w(f) = 1$; Building 0, Chapter 36, P10 says no such $f$ exists. So $R = O$, that is $P = Q$.)
\item[\lab{T3}] (Appears: the points are a move system.) $P \mapsto$ the bag of $P - O$ matches $C(k)$ one-to-one with $\operatorname{Pic}^0(C)$ (T1, T2). Carried back to the points, the law is \[ P \oplus Q = \overline{P * Q}, \qquad O \text{ does nothing}, \qquad \ominus P = \bar P , \] and $P \oplus Q \oplus R = O$ exactly when $P, Q, R$ lie on one line. Associativity of $\oplus$ is P1, carried.
\item[\lab{T4}] (Formulas.) For $P = (x_1, y_1)$, $Q = (x_2, y_2)$ with $x_1 \ne x_2$: $P \oplus Q = (x_3, -y_3)$ with $x_3, y_3$ as in Building 0, Chapter 36, T1. \ $2P$: the same with the tangent's slope $(3x_1^2 + a)/(2y_1)$; $2P = O$ when $y_1 = 0$.
\end{lines}

\subsection*{Propositions on examples}

\begin{lines}
\item[\lab{P5}] (Over $\Q$, a clock.) $y^2 = x^3 + 1$: the points $O$, $(-1, 0)$, $(0, \pm 1)$, $(2, \pm 3)$, and $(2, 3)$ has order $6$: $C(\Q)$ holds a $\Z/6$.
\item[\lab{P6}] (Over $\Q$, no clock.) $y^2 = x^3 - 2$, $P = (3, 5)$: $2P = (\tfrac{129}{100}, -\tfrac{383}{1000})$. The denominators grow with each doubling, and $P$ has no finite order.
\item[\lab{P7}] (Dies when $\Delta = 0$.) On $y^2 = x^3$ T2 fails at the bad point, which every construction through it returns. Leave that point out and the rest is again a move system, matching $(k, +)$ by $(x, y) \mapsto x / y$ (machine-checked over $\Z/p$; classical).
\end{lines}

\subsection*{Remarks}

\begin{lines}
\item[\lab{R1}] (What the method did.) Associativity, the hard line of the usual road, is free (P1): the blur is of a free move system. Existence of a normal form is Euclid-like steps (T1), each a Build. The one ``at most one'' line, T2, is Building 0, Chapter 36, P10, plus one classical step, deferred. A second road to T2 over $\mathbb{C}$ is Chapter III.2.
\item[\lab{R2}] (Kept questions.) The two questions of B2 are all: degree, and ``zeros minus poles of a function''. The same template as Chapter II.4: free thing, gluing, kept questions, normal form.
\item[\lab{R3}] (Old names.) $\Z[C]$: divisors; B1: linear equivalence (lines generate it on a cubic); $\operatorname{Pic}^0(C)$: the degree-zero Picard group; T3: the group law on an elliptic curve; T2: ``no function with a single simple pole'', a case of Riemann--Roch; P4: the chord-and-tangent construction; P6: a point of infinite order; T3 with $k = \Q$ and finitely many generators: Mordell--Weil (classical, reported, not used).
\item[\lab{R4}] (Machine check, \texttt{checks/b1-p3-check.py}.) Over $\Z/p$, $p = 5, 7, 11, 13$, every smooth $(a, b)$: $O$ does nothing, mirror undoes, $\oplus$ commutes, associativity on all triples (on $12$ points per curve when there are more); P5; P6.
\end{lines}

\subsection*{Exercises}

\begin{lines}
\item[\lab{E1}] On $y^2 = x^3 + 1$ over $\Q$, compute $2(2, 3)$ and $3(2, 3)$.
\item[\lab{E2}] Show that $P \oplus P = O$ exactly when $P = O$ or $y(P) = 0$. How many such $P$ are there over $\mathbb{C}$?
\item[\lab{E3}] Why does B1 include the edge line, and what does it glue?
\item[\lab{E4}] Over $\Z/5$, list $C(\Z/5)$ for $y^2 = x^3 + 1$ and find which clock it is.
\end{lines}

\newpage

%% ==================================================================
%%  Chapter III.2  (the torus and the cubic)
%% ==================================================================
\section*{Chapter III.2. The torus carries the law}
\phantomsection\addcontentsline{toc}{section}{\texorpdfstring{Chapter III.2. The torus carries the law \quad $\mathbb{C}/\Lambda$ $\to$ $C(\mathbb{C})$}{Chapter III.2. The torus carries the law}}
\begin{center}
{\normalsize $\mathbb{C}/\Lambda$ $\to$ $C(\mathbb{C})$ \quad$\cdot$\quad from: Chapter III.1; Building 0, Chapters 8, 37 \quad$\cdot$\quad arrow}\\[4pt]
{\small Three points of the torus add to $0$ exactly when their images lie on one line. So the arrow of Building 0, Chapter 37 respects the moves, and $+$ on the torus is the chord law.}
\end{center}

\subsection*{Given}

\begin{lines}
\item[\lab{G1}] (Building 0, Chapter 37.) $\Lambda$, $\mathbb{C}/\Lambda$ with $+$; $\wp$; T4 (zeros and poles equal in number; their sums differ by a member of $\Lambda$); T8 (the arrow $\phi : z \mapsto (\wp(z) : \wp'(z) : 1)$ matches $\mathbb{C}/\Lambda$ with $C(\mathbb{C})$, $C$: $y^2 = 4x^3 - g_2 x - g_3$).
\item[\lab{G2}] (Chapter III.1.) $\oplus$; T3: $P \oplus Q \oplus R = O$ exactly when $P, Q, R$ lie on one line. (The $4$ and the $g_i$ change nothing: scale $y$ by $2$.)
\end{lines}

\subsection*{Theorems}

\begin{lines}
\item[\lab{T1}] (Lines are sums in $\Lambda$.) $\phi(z_1), \phi(z_2), \phi(z_3)$ lie on one line, with counts, exactly when $z_1 + z_2 + z_3 \in \Lambda$. \ (A line $\alpha x + \beta y + \gamma = 0$ with $\beta \ne 0$ gives the elliptic map $\alpha \wp + \beta \wp' + \gamma$, with one triple pole per bag, at $0$. Its zeros are the $z$ whose images lie on the line (G1, T8): three of them, and their sum is in $\Lambda$ (G1, T4). For $\beta = 0$: $\wp - c$ has zeros $z$, $-z$, and the third meeting point is $O = \phi(0)$; again the sum is $0$. Conversely, the line through $\phi(z_1)$ and $\phi(z_2)$ meets $C$ a third time at some $\phi(z_3')$, and $z_3' \sim -z_1 - z_2$.)
\item[\lab{T2}] (Survives: the arrow respects the moves.) $\phi(z_1 + z_2) = \phi(z_1) \oplus \phi(z_2)$. \ (T1 and G2: both sides are the point $R$ with $\phi(z_1) \oplus \phi(z_2) \oplus \ominus R = O$, and $\ominus \phi(z) = \phi(-z)$, the mirror, since $\wp$ is even and $\wp'$ odd.)
\item[\lab{T3}] (A second road to Chapter III.1, T2, over $\mathbb{C}$.) For a cubic reached by T8, $P \mapsto P - O$ is one-to-one: the torus side has no gluing to undo. And associativity of $\oplus$ is the associativity of $+$ on $\mathbb{C}$, carried.
\item[\lab{T4}] (Points of finite order.) The $P$ with $nP = O$ in $C(\mathbb{C})$ match $(\tfrac1n \Lambda)/\Lambda$, so there are $n^2$ of them and they form $\Z/n \times \Z/n$. For $n = 2$: $O$ and the three points with $y = 0$. For $n = 3$: the nine points where the tangent meets $C$ three times.
\end{lines}

\subsection*{Remarks}

\begin{lines}
\item[\lab{R1}] (What the method did.) One arrow. Its respecting $+$ is one line of Building 0, Chapter 37 (T4, the sum of zeros) read through Chapter III.1, T3.
\item[\lab{R2}] (Reported.) Every smooth cubic over $\mathbb{C}$ is reached by some $\Lambda$ (classical, uniformisation). With it, T3 covers every cubic over a $k$ inside $\mathbb{C}$, and the deferred step of Chapter III.1, T2 is no longer needed there. Not used.
\item[\lab{R3}] (Old names.) T1: Abel's theorem on the torus; T2: the addition theorem of $\wp$; T4: $E[n] \cong (\Z/n)^2$; the $n = 3$ points: flexes.
\item[\lab{R4}] (Machine check, \texttt{checks/b1-p3-check.py}.) With $\omega_1 = 1$, $\omega_2 = 0.3 + 1.1i$: for random $z_1, z_2$ and $z_3 = -z_1 - z_2 + \omega_1 + \omega_2$ the three images have $\det$ below $10^{-6}$ (relative); moving $z_3$ by $0.13$ gives $\det$ above $3 \cdot 10^{-3}$.
\end{lines}

\subsection*{Exercises}

\begin{lines}
\item[\lab{E1}] Use T2 to write $\wp(2z)$ in terms of $\wp(z)$ and $\wp'(z)$.
\item[\lab{E2}] Why are the points with $nP = O$ over $\R$ fewer than $n^2$?
\item[\lab{E3}] Show from T1 that $\phi(z)$ is a point where the tangent meets $C$ three times exactly when $3z \in \Lambda$.
\end{lines}

\newpage

%% ==================================================================
%%  Chapter III.3  (points over Z/p)
%% ==================================================================
\section*{Chapter III.3. Counting points over a clock}
\phantomsection\addcontentsline{toc}{section}{\texorpdfstring{Chapter III.3. Counting points over a clock \quad cubics over $\Z/p$ $\to$ their sizes}{Chapter III.3. Counting points over a clock}}
\begin{center}
{\normalsize cubics over $\Z/p$ $\to$ their sizes \quad$\cdot$\quad from: Chapters II.2, II.4, III.1; Building 0, Chapters 3, 17 \quad$\cdot$\quad count}\\[4pt]
{\small Over $\Z/p$ the move system of Chapter III.1 is finite. Its size is $p + 1$ plus a sum of $p$ signs, and the signs nearly cancel.}
\end{center}

\subsection*{Given}

\begin{lines}
\item[\lab{G1}] (Chapter III.1.) $C(\Z/p)$ with $\oplus$, for $p \ne 2, 3$ and $\Delta \ne 0$ in $\Z/p$.
\item[\lab{G2}] (Chapter II.2, T1, T2.) Lagrange; the order of a move divides the size.
\item[\lab{G3}] (Chapter II.4, T2.) A finite $\mathsf{comm}$ move system is $\Z/e_1 \times \dots \times \Z/e_r$, $e_1 \mid \dots \mid e_r$, in one way.
\item[\lab{G4}] (Building 0, Chapter 3, T4, T5.) $\Z/p$ is a field; $c^{p} = c$.
\end{lines}

\subsection*{Definitions}

\begin{lines}
\item[\lab{D1}] (The square sign.) For $c$ in $\Z/p$: $\chi(c) := 0$ if $c = 0$; $1$ if $c = y^2$ for some $y$; $-1$ otherwise. \ $N_p := |C(\Z/p)|$, $O$ included.
\end{lines}

\subsection*{Propositions}

\begin{lines}
\item[\lab{P1}] (How many square roots.) $y^2 = c$ has $1 + \chi(c)$ answers in $\Z/p$. \ ($y$ and $-y$ differ unless $y = 0$, as $2 \ne 0$; at most two by Building 0, Chapter 7, T4 over $\Z/p$.)
\item[\lab{P2}] (The count.) \[ N_p \;=\; p + 1 + \sum_{x \in \Z/p} \chi\big(x^3 + ax + b\big). \] (P1 for each $x$, and $O$.)
\item[\lab{P3}] (Euler's test.) $\chi(c) = c^{(p-1)/2}$ in $\Z/p$, read as $0$, $1$, $-1$. \ (The squares $c \neq 0$ are $(p - 1)/2$ in number (P1) and each has $c^{(p-1)/2} = y^{p-1} = 1$ (G4); $x^{(p-1)/2} - 1$ has no more roots than that.)
\end{lines}

\subsection*{Theorems}

\begin{lines}
\item[\lab{T1}] (Lagrange on the curve.) For every $P$ in $C(\Z/p)$, $N_p\,P = O$. \ (G2.)
\item[\lab{T2}] (At most two clocks.) $C(\Z/p)$ matches $\Z/e_1 \times \Z/e_2$ with $e_1 \mid e_2$ (G3; at most two factors because the $P$ with $\ell P = O$, $\ell$ prime, number at most $\ell^2$: classical for $\Z/p$; over $\mathbb{C}$ it is Chapter III.2, T4). Also $e_1 \mid p - 1$ (classical, reported).
\item[\lab{T3}] (The signs nearly cancel. Reported; its road is Building 2.) \[ |N_p - p - 1| \;\le\; 2\sqrt{p} . \] The sum in P2 has $p$ terms $\pm 1$ and $0$; if they were independent coin tosses its size would be about $\sqrt p$ (Building 0, Chapter 17's scale). T3 says the curve's signs never do worse than twice that. \\
The road: the counts $N_{p^m}$ over the floors with $p^m$ members form a ladder; one rough bound $|N_{p^m} - p^m - 1| \le c\,p^{m/2}$, the same $c$ on every rung, is amplified along the ladder to the constant $2$ (Amplify, Building 2). Two inputs of that road are classical and reported: the shape $N_{p^m} = p^m + 1 - \alpha^m - \beta^m$ with $\alpha\beta = p$, and the rough bound.
\end{lines}

\subsection*{Remarks}

\begin{lines}
\item[\lab{R1}] (What the method did.) Counting, and two lines from Part II. The real content, T3, is a cancellation line: the edge the survey of the book keeps meeting (Building 0, Chapter 17, P6).
\item[\lab{R2}] (Old names.) $\chi$: the Legendre symbol; P3: Euler's criterion; T3: Hasse's theorem (the Riemann hypothesis for elliptic curves over finite fields); $\alpha, \beta$: Frobenius eigenvalues; $p + 1 - N_p$: the trace of Frobenius.
\item[\lab{R3}] (Machine check, \texttt{checks/b1-p3-check.py}.) For every prime $5 \le p < 60$ and every smooth $(a, b)$: P2 against direct listing; T1 (each point's order divides $N_p$); T2 ($e_1 e_2 = N_p$, $e_1 \mid e_2$, $e_1 \mid p - 1$); T3, with the worst case $|N_p - p - 1| / 2\sqrt p = 0.991$.
\end{lines}

\subsection*{Exercises}

\begin{lines}
\item[\lab{E1}] Count $C(\Z/7)$ for $y^2 = x^3 + 2$ by P2, and check T3.
\item[\lab{E2}] Show that for $y^2 = x^3 + x$ and $p \equiv 3 \pmod 4$, $N_p = p + 1$. (Pair $x$ with $-x$, and use $\chi(-1) = -1$.)
\item[\lab{E3}] Use T1 to find the order of $(0, 1)$ on $y^2 = x^3 + 1$ over $\Z/5$ without adding.
\end{lines}

\end{document}
